\documentclass[UTF8,12pt]{ctexart}
\usepackage[a4paper,margin=24mm]{geometry}
\usepackage{amsmath,amssymb,bm,graphicx,adjustbox}
\newcommand{\fitmath}[1]{\begin{adjustbox}{max width=\linewidth}$\displaystyle #1$\end{adjustbox}}
\setlength{\parindent}{0pt}
\setlength{\parskip}{.7em}
\sloppy
\allowdisplaybreaks
\begin{document}
\section*{2018 年考研数学三 · 真题}
\subsection*{原 PDF 第 6 页}

\subsection*{2018年全国硕士研究生招生考试试题}

\subsection*{一、选择题（本题共8小题，每小题4分，共32分。在每小题给出的四个选项中，只有一项符合题目要求，把所选项前的字母填在题后的括号内。）}

（1）下列函数中，在 \(\displaystyle x=0\) 处不可导的是（　）。


（A）\(\displaystyle f(x)=|x|\sin|x|\)。 （B）\(\displaystyle f(x)=|x|\sin\sqrt{|x|}\)。 （C）\(\displaystyle f(x)=\cos|x|\)。 （D）\(\displaystyle f(x)=\cos\sqrt{|x|}\)。


（2）设函数 \(\displaystyle f(x)\) 在 \(\displaystyle [0,\allowbreak 1]\) 上二阶可导，且 \(\displaystyle \int_0^1f(x)\,\mathrm dx=0\)，则（　）。


（A）当 \(\displaystyle f^{\prime}(x)<0\) 时，\(\displaystyle f(\dfrac{\displaystyle 1}{\displaystyle 2})<0\)。 （B）当 \(\displaystyle f^{\prime\prime}(x)<0\) 时，\(\displaystyle f(\dfrac{\displaystyle 1}{\displaystyle 2})<0\)。 （C）当 \(\displaystyle f^{\prime}(x)>0\) 时，\(\displaystyle f(\dfrac{\displaystyle 1}{\displaystyle 2})<0\)。 （D）当 \(\displaystyle f^{\prime\prime}(x)>0\) 时，\(\displaystyle f(\dfrac{\displaystyle 1}{\displaystyle 2})<0\)。


（3）设 \(\displaystyle M=\displaystyle \int_{-\pi/2}^{\pi/2}\dfrac{\displaystyle (1+x)^2}{\displaystyle 1+x^2}\,\mathrm dx\)，\(\displaystyle N=\displaystyle \int_{-\pi/2}^{\pi/2}\dfrac{\displaystyle 1+x}{\displaystyle e^x}\,\mathrm dx\)，\(\displaystyle K=\displaystyle \int_{-\pi/2}^{\pi/2}(1+\sqrt{\cos x})\,\mathrm dx\)，则（　）。


（A）\(\displaystyle M>N>K\)。 （B）\(\displaystyle M>K>N\)。 （C）\(\displaystyle K>M>N\)。 （D）\(\displaystyle K>N>M\)。


（4）设某产品的成本函数 \(\displaystyle C(Q)\) 可导，其中 \(\displaystyle Q\) 为产量。若产量为 \(\displaystyle Q_0\) 时平均成本最小，则（　）。


（A）\(\displaystyle C^{\prime}(Q_0)=0\)。 （B）\(\displaystyle C^{\prime}(Q_0)=C(Q_0)\)。 （C）\(\displaystyle C^{\prime}(Q_0)=Q_0C(Q_0)\)。 （D）\(\displaystyle Q_0C^{\prime}(Q_0)=C(Q_0)\)。


（5）下列矩阵中，与矩阵 \(\displaystyle \begin{pmatrix}1&1&0\\0&1&1\\0&0&1\end{pmatrix}\) 相似的为（　）。


（A）\(\displaystyle \begin{pmatrix}1&1&-1\\0&1&1\\0&0&1\end{pmatrix}\)。 （B）\(\displaystyle \begin{pmatrix}1&0&-1\\0&1&1\\0&0&1\end{pmatrix}\)。 （C）\(\displaystyle \begin{pmatrix}1&1&-1\\0&1&0\\0&0&1\end{pmatrix}\)。 （D）\(\displaystyle \begin{pmatrix}1&0&-1\\0&1&0\\0&0&1\end{pmatrix}\)。


（6）设 \(\displaystyle A,\allowbreak B\) 为 \(\displaystyle n\) 阶矩阵，记 \(\displaystyle r(X)\) 为矩阵 \(\displaystyle X\) 的秩，\(\displaystyle (X,\allowbreak Y)\) 表示分块矩阵，则（　）。


（A）\(\displaystyle r(A,\allowbreak AB)=r(A)\)。 （B）\(\displaystyle r(A,\allowbreak BA)=r(A)\)。 （C）\(\displaystyle r(A,\allowbreak B)=\displaystyle \max\{r(A),\allowbreak r(B)\}\)。 （D）\(\displaystyle r(A,\allowbreak B)=r(A^{\mathrm T},\allowbreak B^{\mathrm T})\)。


（7）设随机变量 \(\displaystyle X\) 的概率密度 \(\displaystyle f(x)\) 满足 \(\displaystyle f(1+x)=f(1-x)\)，且 \(\displaystyle \int_0^2f(x)\,\mathrm dx=0.6\)，则 \(\displaystyle P\{X<0\}=\)（　）。


（A）0.2。 （B）0.3。 （C）0.4。 （D）0.5。


（8）设 \(\displaystyle X_1,\allowbreak X_2,\allowbreak \cdots,\allowbreak X_n\;(n\ge2)\) 为来自总体 \(\displaystyle N(\mu,\allowbreak \sigma^2)\;(\sigma>0)\) 的简单随机样本。令 \(\displaystyle \bar X=\dfrac{\displaystyle 1}{\displaystyle n}\displaystyle \sum_{i=1}^nX_i\)，\(\displaystyle S=\sqrt{\dfrac{\displaystyle 1}{\displaystyle n-1}\displaystyle \sum_{i=1}^n(X_i-\bar X)^2}\)，\(\displaystyle S^*=\sqrt{\dfrac{\displaystyle 1}{\displaystyle n}\displaystyle \sum_{i=1}^n(X_i-\mu)^2}\)，则（　）。


（A）\(\displaystyle \dfrac{\displaystyle \sqrt n(\bar X-\mu)}{\displaystyle S}\sim t(n)\)。 （B）\(\displaystyle \dfrac{\displaystyle \sqrt n(\bar X-\mu)}{\displaystyle S}\sim t(n-1)\)。


\clearpage

\subsection*{原 PDF 第 7 页}

（8）（续）（C）\(\displaystyle \dfrac{\displaystyle \sqrt n(\bar X-\mu)}{\displaystyle S^*}\sim t(n)\)。 （D）\(\displaystyle \dfrac{\displaystyle \sqrt n(\bar X-\mu)}{\displaystyle S^*}\sim t(n-1)\)。


\subsection*{二、填空题（本题共6小题，每小题4分，共24分，把答案填在题中横线上。）}

（9）曲线 \(\displaystyle y=x^2+2\ln x\) 在其拐点处的切线方程是\_\_\_\_\_\_。


（10）\(\displaystyle \int e^x\arcsin\sqrt{1-e^{2x}}\,\mathrm dx=\)\_\_\_\_\_\_。


（11）差分方程 \(\displaystyle \Delta^2y_x-y_x=5\) 的通解为\_\_\_\_\_\_。


（12）设函数 \(\displaystyle f(x)\) 满足 \(\displaystyle f(x+\Delta x)-f(x)=2xf(x)\Delta x+o(\Delta x)\;(\Delta x\to0)\)，且 \(\displaystyle f(0)=2\)，则 \(\displaystyle f(1)=\)\_\_\_\_\_\_。


（13）设 \(\displaystyle A\) 为3阶矩阵，\(\displaystyle \alpha_1,\allowbreak \alpha_2,\allowbreak \alpha_3\) 是线性无关的向量组。若 \(\displaystyle A\alpha_1=\alpha_1+\alpha_2\)，\(\displaystyle A\alpha_2=\alpha_2+\alpha_3\)，\(\displaystyle A\alpha_3=\alpha_1+\alpha_3\)，则 \(\displaystyle |A|=\)\_\_\_\_\_\_。


（14）随机事件 \(\displaystyle A,\allowbreak B,\allowbreak C\) 相互独立，且 \(\displaystyle P(A)=P(B)=P(C)=\dfrac{\displaystyle 1}{\displaystyle 2}\)，则 \(\displaystyle P(AC\mid A\cup B)=\)\_\_\_\_\_\_。


\subsection*{三、解答题（本题共9小题，共94分，解答应写出文字说明、证明过程或演算步骤。）}

（15）（本题满分10分）已知实数 \(\displaystyle a,\allowbreak b\) 满足 \(\displaystyle \lim_{x\to+\infty}[(ax+b)e^{1/x}-x]=2\)，求 \(\displaystyle a,\allowbreak b\)。


（16）（本题满分10分）设平面区域 \(\displaystyle D\) 由曲线 \(\displaystyle y=\sqrt{3(1-x^2)}\) 与直线 \(\displaystyle y=\sqrt3x\) 及 \(\displaystyle y\) 轴围成。计算二重积分 \(\displaystyle \iint_D x^2\,\mathrm dx\mathrm dy\)。


\clearpage

\subsection*{原 PDF 第 8 页}

（17）（本题满分10分）将长为2m的铁丝分成三段，依次围成圆、正方形与正三角形。三个图形的面积之和是否存在最小值？若存在，求出最小值。


（18）（本题满分10分）已知 \(\displaystyle \cos2x-\dfrac{\displaystyle 1}{\displaystyle (1+x)^2}=\displaystyle \sum_{n=0}^{\infty}a_nx^n\;(-1<x<1)\)，求 \(\displaystyle a_n\)。


（19）（本题满分10分）设数列 \(\displaystyle \{x_n\}\) 满足：\(\displaystyle x_1>0\)，\(\displaystyle x_ne^{x_{n+1}}=e^{x_n}-1\;(n=1,\allowbreak 2,\allowbreak \cdots)\)。证明 \(\displaystyle \{x_n\}\) 收敛，并求 \(\displaystyle \lim_{n\to\infty}x_n\)。


（20）（本题满分11分）设实二次型 \(\displaystyle f(x_1,\allowbreak x_2,\allowbreak x_3)=(x_1-x_2+x_3)^2+(x_2+x_3)^2+(x_1+ax_3)^2\)，其中 \(\displaystyle a\) 是参数。（I）求 \(\displaystyle f(x_1,\allowbreak x_2,\allowbreak x_3)=0\) 的解；（II）求 \(\displaystyle f(x_1,\allowbreak x_2,\allowbreak x_3)\) 的规范形。


\clearpage

\subsection*{原 PDF 第 9 页}

（21）（本题满分11分）已知 \(\displaystyle a\) 是常数，且矩阵 \(\displaystyle A=\begin{pmatrix}1&2&a\\1&3&0\\2&7&-a\end{pmatrix}\) 可经初等列变换化为矩阵 \(\displaystyle B=\begin{pmatrix}1&a&2\\0&1&1\\-1&1&1\end{pmatrix}\)。（I）求 \(\displaystyle a\)；（II）求满足 \(\displaystyle AP=B\) 的可逆矩阵 \(\displaystyle P\)。


（22）（本题满分11分）设随机变量 \(\displaystyle X\) 与 \(\displaystyle Y\) 相互独立，\(\displaystyle X\) 的概率分布为 \(\displaystyle P\{X=1\}=P\{X=-1\}=\dfrac{\displaystyle 1}{\displaystyle 2}\)，\(\displaystyle Y\) 服从参数为 \(\displaystyle \lambda\) 的泊松分布。令 \(\displaystyle Z=XY\)。（I）求 \(\displaystyle \operatorname{Cov}(X,\allowbreak Z)\)；（II）求 \(\displaystyle Z\) 的概率分布。


（23）（本题满分11分）设总体 \(\displaystyle X\) 的概率密度为 \(\displaystyle f(x;\sigma)=\dfrac{\displaystyle 1}{\displaystyle 2\sigma}e^{-|x|/\sigma}\)，\(\displaystyle -\infty<x<+\infty\)，其中 \(\displaystyle \sigma\in(0,\allowbreak +\infty)\) 为未知参数，\(\displaystyle X_1,\allowbreak X_2,\allowbreak \cdots,\allowbreak X_n\) 为来自总体 \(\displaystyle X\) 的简单随机样本。记 \(\displaystyle \sigma\) 的最大似然估计量为 \(\displaystyle \hat\sigma\)。（I）求 \(\displaystyle \hat\sigma\)；（II）求 \(\displaystyle E(\hat\sigma),\allowbreak D(\hat\sigma)\)。


\clearpage
\end{document}
